%% ****** Start of file apstemplate.tex ****** %
%%
%%
%%   This file is part of the APS files in the REVTeX 4 distribution.
%%   Version 4.1r of REVTeX, August 2010
%%
%%
%%   Copyright (c) 2001, 2009, 2010 The American Physical Society.
%%
%%   See the REVTeX 4 README file for restrictions and more information.
%%
%
% This is a template for producing manuscripts for use with REVTEX 4.0
% Copy this file to another name and then work on that file.
% That way, you always have this original template file to use.
%
% Group addresses by affiliation; use superscriptaddress for long
% author lists, or if there are many overlapping affiliations.
% For Phys. Rev. appearance, change preprint to twocolumn.
% Choose pra, prb, prc, prd, pre, prl, prstab, prstper, or rmp for journal
%  Add 'draft' option to mark overfull boxes with black boxes
%  Add 'showpacs' option to make PACS codes appear
%  Add 'showkeys' option to make keywords appear
%\documentclass[aps,prl,preprint,superscriptaddress,twocolumn]{revtex4-1}
\documentclass[aps,prd,twocolumn,groupedaddress]{revtex4-2}
%\documentclass[aps,prl,reprint,groupedaddress,showpacs]{revtex4-1}
\pdfoutput=1

               			
\usepackage{amssymb}
\usepackage{amsmath}
\usepackage{graphicx}
\usepackage{enumerate}
\usepackage{mathtools}
\usepackage{color}
\usepackage{bbold}
\usepackage{subfigure}
\usepackage{slashed}
\usepackage{xcolor}
\usepackage{bm}
\usepackage{multirow}
\usepackage{pifont}


% You should use BibTeX and apsrev.bst for references
% Choosing a journal automatically selects the correct APS
% BibTeX style file (bst file), so only uncomment the line
% below if necessary.
%\bibliographystyle{apsrev4-1}

\begin{document}

% Use the \preprint command to place your local institutional report
% number in the upper righthand corner of the title page in preprint mode.
% Multiple \preprint commands are allowed.
% Use the 'preprintnumbers' class option to override journal defaults
% to display numbers if necessary

%Title of paper
\title{Supplemental material: Thermodynamics of oscillating neutrinos}

% repeat the \author .. \affiliation  etc. as needed
% \email, \thanks, \homepage, \altaffiliation all apply to the current
% author. Explanatory text should go in the []'s, actual e-mail
% address or url should go in the {}'s for \email and \homepage.
% Please use the appropriate macro foreach each type of information

% \affiliation command applies to all authors since the last
% \affiliation command. The \affiliation command should follow the
% other information
% \affiliation can be followed by \email, \homepage, \thanks as well.
\author{Lucas Johns}
%\email[]{NASA Einstein Fellow (ljohns@berkeley.edu)}
\affiliation{Departments of Astronomy \& Physics, University of California, Berkeley, CA 94720, USA}








\maketitle

\renewcommand{\theequation}{S.\arabic{equation}}





\section{Kinematic decoherence}

In the main text we report that the maximum-entropy state of a neutrino traveling in vacuum is
\begin{equation}
\rho^\textrm{eq}_{ij} = \rho^\textrm{in}_{ii} \delta_{ij}, \label{eq:vacanswer}
\end{equation}
where $i$ and $j$ are mass indices and $\rho^\textrm{in}$ is the flavor/mass distribution at the source. The calculation is simple but illustrative.

We use $\overline{\rho}$ to indicate the coarse-grained density matrix $\rho$. The particular coarse-graining we employ is inessential. It could be an average over a window of energies, for example, which would correspond to an experiment in which we make finite-resolution detections of neutrinos in a beam of average energy $p$. We will simply assume that there is \textit{some} coarse-graining---some conduit by which information about the system is lost---that justifies the maximization of entropy.

Because the neutrinos are in vacuum, the Hamiltonian is
\begin{equation}
H = p + M^2 / 2p.
\end{equation}
The conserved quantities in this scenario are $\textrm{Tr}\left[ H^n \overline{\rho} \right]$ for $0 \leq n \leq N_f$, with the $n > 1$ invariants following from the time-independence of the Hamiltonian. As is true generally, the nonlinear ($n > 1$) quantities $\textrm{Tr}[\overline{\rho}^n]$ are not conserved because $\overline{\rho^n} \neq \overline{\rho}^n$. The statistical ensemble therefore has expectation values
\begin{align}
&\textrm{Tr} \left[ \overline{\rho} \right] = \overline{\textrm{Tr}\left[ \rho \right]} \equiv \mathcal{N}, \notag \\
&\textrm{Tr}\left[ H^n \overline{\rho} \right] = \overline{\textrm{Tr}\left[ H^n \rho \right]} \equiv \mathcal{U}_n, \label{eq:vacconslaws}
\end{align}
where the lower line is for $n \geq 1$.

Let $\eta$ and $\lbrace \beta_n \rbrace$ be the Lagrange multipliers associated with $\mathcal{N}$ and $\lbrace \mathcal{U}_n \rbrace$, respectively. Then the entropy with constraints is
\begin{align}
S' = &- \textrm{Tr} \left[ \overline{\rho} \log \overline{\rho} + \left( 1 - \overline{\rho} \right) \log \left( 1 - \overline{\rho} \right) \right] \notag \\
&- \eta \left( \textrm{Tr} \left[ \overline{\rho} \right] - \mathcal{N} \right) + \sum_n \beta_n \left( \textrm{Tr} \left[ H^n \overline{\rho} \right] - \mathcal{U}_n \right).
\end{align}
From
\begin{equation}
\frac{\delta S'}{\delta \overline{\rho}} \bigg|_{\rho^\textrm{eq}} = 0
\end{equation}
we obtain
\begin{equation}
\rho^\textrm{eq} = \frac{1}{\exp\left[ \sum_{n=1}^{N_f} \beta_n \left( p + \frac{M^2}{2p} \right)^n  - \eta\right] + 1}. \label{eq:vaceq}
\end{equation}
The distribution is diagonal in the mass basis. Working in this basis, the conservation laws in Eq.~\eqref{eq:vacconslaws} give
\begin{equation}
\sum_{i=1}^{N_f} \left( p + \frac{m_i^2}{2p} \right)^n \rho_{ii}^\textrm{eq} = \sum_{i=1}^{N_f} \left( p + \frac{m_i^2}{2p} \right)^n \rho_{ii}^\textrm{in}
\end{equation}
for $0 \leq n \leq N_f$. Hence the problem is solved by Eq.~\eqref{eq:vacanswer} as claimed.

This example is simple enough that we can find the solution without knowing $\lbrace \eta, \beta_n \rbrace$, but an explicit determination of the $N_f + 1$ Lagrange multipliers could be carried out numerically if desired. To do that, we would solve the $N_f + 1$ self-consistency conditions
\begin{align}
&\mathcal{N} = \textrm{Tr} \left[ \rho^\textrm{eq} \left( \eta, \beta_n \right) \right], \notag \\
&\mathcal{U}_n = \textrm{Tr} \left[ H^n \rho^\textrm{eq} \left( \eta, \beta_n \right) \right],
\end{align}
which result from mandating agreement between the conservation laws in Eq.~\eqref{eq:vacconslaws} and the functional form of $\rho^\textrm{eq}$ in Eq.~\eqref{eq:vaceq}.

Incidentally, a similar calculation shows up in the thermodynamics of massless neutrinos that are isolated from their surroundings but interact with each other through $2 \rightarrow 2$ collisions. For a single neutrino species, the conserved quantities would be temperature $T$ and chemical potential $\mu$, and they would have to be obtained by solving the self-consistency relations among number density $n$, energy density $u$, and the Fermi--Dirac distribution $f_\textrm{FD} (p, T, \mu)$:
\begin{align}
&n = \int \frac{d^3 \boldsymbol{p}}{(2\pi)^3} f_\textrm{FD} \left( p, T, \mu \right), \notag \\
&u = \int \frac{d^3 \boldsymbol{p}}{(2\pi)^3} p f_\textrm{FD} \left( p, T, \mu \right).
\end{align} 
These types of calculations, where one must infer intensive thermodynamic parameters from known extensive ones, arise in considering the thermalization of isolated systems because such systems act as their own environments.













\section{The miscidynamic equation}

In this section we derive the miscidynamic equation
\begin{equation}
i \left( \partial_t + \boldsymbol{\hat{p}} \cdot \partial_{\boldsymbol{x}} \right) \rho_\nu^\textrm{eq} (t, \boldsymbol{x}, \boldsymbol{p}) = iC_\nu^\textrm{eq} (t, \boldsymbol{x}, \boldsymbol{p}) \label{eq:miscmisc}
\end{equation}
by coarse-graining the neutrino kinetic equation
\begin{align}
i \left( \partial_t + \boldsymbol{\hat{p}} \cdot \partial_{\boldsymbol{x}} \right) \rho_\nu \left( t, \boldsymbol{x}, \boldsymbol{p} \right) = &\left[ H_\nu (t, \boldsymbol{x}, \boldsymbol{p} ), \rho_\nu \left( t, \boldsymbol{x}, \boldsymbol{p} \right) \right] \notag \\
&~~~~~~~~ + iC_\nu (t, \boldsymbol{x}, \boldsymbol{p}). \label{eq:miscqke}
\end{align}

First we define the coarse-graining operator 
\begin{equation}
\left\langle \cdot \right\rangle (T, \boldsymbol{X}, \boldsymbol{p}) \equiv \frac{1}{V \Delta t} \int_{\mathcal{R}_{\boldsymbol{x}}} d^3\boldsymbol{x}' \int_t^{t+\Delta t} dt' \left( \cdot \right) \left( t', \boldsymbol{x}', \boldsymbol{p} \right), \label{eq:cgop}
\end{equation}
which averages over a region $\mathcal{R}_{\boldsymbol{x}}$ with volume $V$ centered at location $\boldsymbol{x}$ and over a time step $\Delta t$. We are not coarse-graining momentum $\boldsymbol{p}$, though nothing prevents us from doing so. The use of $T = t$ and $\boldsymbol{X} = \boldsymbol{x}$ on the left-hand side emphasizes that these are the coordinates of coarse-grained variables. (Elsewhere in the paper we use $T$ to denote temperature. In this section only we use it as a time coordinate.) Thinking of the coordinates themselves as being coarse-grained is justified by the idempotent property
\begin{equation}
\left\langle \left\langle \cdot \right\rangle \right\rangle \cong \left\langle \cdot \right\rangle \label{eq:idem}
\end{equation}
of the coarse-graining operator. We can therefore regard $\langle \cdot \rangle$ as acting on the fine-grained $t$ and $\boldsymbol{x}$ but seeing the coarse-grained $T$ and $\boldsymbol{X}$ as constants. Eq.~\eqref{eq:idem} holds as long as we take $\Delta t$ and $V$ to be small relative to the scales of global variations (but still large enough that small-scale features are smoothed out, otherwise we gain nothing from the operator). In the application to $\rho_\nu$, we are taking $\Delta t \sim t_\textrm{box}$ and $V \sim l_\textrm{box}^3$ and using the scale separations $t_\textrm{osc} \ll t_\textrm{astro}$ and $l_\textrm{osc} \ll l_\textrm{astro}$, as described in the main text.

Our second step is to write
\begin{equation}
\rho_\nu (t, \boldsymbol{x}, \boldsymbol{p}) = \left\langle \rho_\nu \right\rangle (T, \boldsymbol{X}, \boldsymbol{p})  + \delta \rho_\nu (t, \boldsymbol{x}, \boldsymbol{p}), \label{eq:miscexpand}
\end{equation}
which defines the deviation $\delta \rho_\nu$ from the average. We have
\begin{equation}
\left\langle \delta \rho_\nu \right\rangle \cong 0, \label{eq:miscavgdelta}
\end{equation}
consistent with Eq.~\eqref{eq:idem}.

Now we use Eq.~\eqref{eq:miscexpand} in Eq.~\eqref{eq:miscqke} and act on the resulting equation with $\left\langle \cdot \right\rangle$. Terms with no factors of $\delta \rho_\nu$ are unchanged according to Eq.~\eqref{eq:idem}. Terms that are linear in $\delta \rho_\nu$ vanish by Eq.~\eqref{eq:miscavgdelta}. The remaining terms contain correlators like $\left\langle \delta \rho_\nu \delta \rho_\nu \right\rangle$, $\left\langle \delta \rho_\nu \delta \rho_{\bar{\nu}} \right\rangle$, and so on, which arise from the self-interactions in $[ H_\nu, \rho_\nu ]$ and $i C_\nu$. All $\delta \rho_\nu$ and $\delta \rho_{\bar{\nu}}$ factors are at the same $t$ and $\boldsymbol{x}$ but can have different $\boldsymbol{p}$.

By the ergodic hypothesis,
\begin{equation}
\left\langle \rho_\nu \right\rangle (T, \boldsymbol{X}, \boldsymbol{p}) \cong \overline{\rho_{\nu,\boldsymbol{p}}} (T, \boldsymbol{X}). \label{eq:miscavgerg}
\end{equation}
The left-hand expression is the time- and position-averaged dynamical solution at $(T, \boldsymbol{X})$, as defined in Eq.~\eqref{eq:cgop}. The right-hand expression is the ensemble average of possible microstates of the $iC_\nu = 0 $ system with parameters (\textit{e.g.}, electron density) set to the values at $(T, \boldsymbol{X})$. The overline notation matches the main text.

The implication of Eq.~\eqref{eq:miscavgerg} is that $\left\langle \delta \rho_\nu \delta \rho_\nu \right\rangle$ and other such correlators describe the statistics of fluctuations around the ensemble average. All of these terms go to zero in the thermodynamic limit. After switching to ensemble averages and dropping fluctuation correlations, we are left with 
\begin{equation}
i \left( \partial_T + \boldsymbol{\hat{p}} \cdot \partial_{\boldsymbol{X}} \right) \overline{\rho_{\nu,\boldsymbol{p}}} = \left[ H_{\nu,\boldsymbol{p}} \left[ \overline{\rho_\nu} \right], \overline{\rho_{\nu,\boldsymbol{p}}} \right] + iC_\nu \left[ \overline{\rho_\nu} \right].
\end{equation}
All functions in this equation depend on $(T, \boldsymbol{X})$. We omit the $\boldsymbol{p}$ subscripts in the arguments of $H_{\nu,\boldsymbol{p}}$ and $iC_\nu$ to indicate that these are functionals of the coarse-grained density matrices at all momenta.

The last step is to identify
\begin{equation}
\overline{\rho_{\nu,\boldsymbol{p}}} (T, \boldsymbol{X}) = \rho_\nu^\textrm{eq} (T, \boldsymbol{X}, \boldsymbol{p}),
\end{equation}
where $\rho_\nu^\textrm{eq}$ refers to the equilibrium distribution obtained in the main text. The coordinates specify when and where the parameters are drawn that enter into the Hamiltonian and determine the values of the conserved quantities. Observing that $\rho_\nu^\textrm{eq}$ commutes with its Hamiltonian and reverting to symbols $\boldsymbol{t}$ and $\boldsymbol{x}$, we arrive at Eq.~\eqref{eq:miscmisc}.
















\bibliography{all_papers}

\end{document}